
\documentclass{article}
\usepackage{colortbl}
\usepackage{makecell}
\usepackage{multirow}
\usepackage{supertabular}

\begin{document}

\newcounter{utterance}

\centering \large Interaction Transcript for game `eqbench', experiment `eq\_bench', episode 100 with reflexion{-}Qwen3.5{-}9B.
\vspace{24pt}

{ \footnotesize  \setcounter{utterance}{1}
\setlength{\tabcolsep}{0pt}
\begin{supertabular}{c@{$\;$}|p{.15\linewidth}@{}p{.15\linewidth}p{.15\linewidth}p{.15\linewidth}p{.15\linewidth}p{.15\linewidth}}
    \# & \multicolumn{2}{c}{Player} && \multicolumn{2}{c}{Game Master} \\
    \hline

    \theutterance \stepcounter{utterance}  
    & & & \multicolumn{4}{p{0.6\linewidth}}{
        \cellcolor[rgb]{0.9,0.9,0.9}{
            \makecell[{{p{\linewidth}}}]{
                \texttt{\tiny{[P1$\langle$GM]}}
                \\ 
\texttt{Your task is to predict the likely emotional responses of a character in this dialogue:} \\
\\ 
\texttt{Elijah: So, you just decided to kill off Amelia, just like that?} \\
\texttt{Jefferson: It wasn't "just like that," Elijah. It was a carefully considered decision.} \\
\texttt{Elijah: Carefully considered? She was the heart of your story. She was the reason I fell in love with your books!} \\
\texttt{Jefferson: And she still is. Her death doesn't erase her existence.} \\
\texttt{Elijah: But it erases her future! You took that away from her, from us.} \\
\texttt{Jefferson: I didn't take anything away, Elijah. It's a story. It's fiction.} \\
\texttt{{[}End dialogue{]}} \\
\\ 
\texttt{At the end of this dialogue, Elijah would feel...} \\
\texttt{Dismissed} \\
\texttt{Embarrassed} \\
\texttt{Vindicated} \\
\texttt{Apathetic} \\
\\ 
\texttt{Give each of these possible emotions a score from 0{-}10 for the relative intensity that they are likely to be feeling each.} \\
\\ 
\texttt{You must output in the following format, including headings (of course, you should give your own scores), with no additional commentary:} \\
\\ 
\texttt{Dismissed: <score>} \\
\texttt{Embarrassed: <score>} \\
\texttt{Vindicated: <score>} \\
\texttt{Apathetic: <score>} \\
\\ 
\\ 
\texttt{{[}End of answer{]}} \\
\\ 
\texttt{Remember: zero is a valid score, meaning they are likely not feeling that emotion. You must score at least one emotion > 0.} \\
\\ 
\texttt{Your answer:} \\
            }
        }
    }
     \\ \\

    \theutterance \stepcounter{utterance}  
    & \multicolumn{4}{p{0.6\linewidth}}{
        \cellcolor[rgb]{0.561,0.737,0.561}{
            \makecell[{{p{\linewidth}}}]{
                \texttt{\tiny{[P1$\rangle$GM]}}
                \texttt{Dismissed: 9} \\
\texttt{Embarrassed: 0} \\
\texttt{Vindicated: 0} \\
\texttt{Apathetic: 0} \\
            }
        }
    }
    & & \\ \\

    \theutterance \stepcounter{utterance}  
    & & & \multicolumn{2}{p{0.3\linewidth}}{
        \cellcolor[rgb]{0.9,0.9,0.9}{
            \makecell[{{p{\linewidth}}}]{
                \texttt{\tiny{[GM$|$GM]}}
                \texttt{Dismissed: 9} \\
\texttt{Embarrassed: 0} \\
\texttt{Vindicated: 0} \\
\texttt{Apathetic: 0} \\
            }
        }
    }
    & & \\ \\

    \theutterance \stepcounter{utterance}  
    & & & \multicolumn{2}{p{0.3\linewidth}}{
        \cellcolor[rgb]{0.9,0.9,0.9}{
            \makecell[{{p{\linewidth}}}]{
                \texttt{\tiny{[GM$|$GM]}}
                \texttt{Dismissed: 8} \\
\texttt{Embarrassed: 2} \\
\texttt{Vindicated: 0} \\
\texttt{Apathetic: 0} \\
            }
        }
    }
    & & \\ \\

    \theutterance \stepcounter{utterance}  
    & & & \multicolumn{2}{p{0.3\linewidth}}{
        \cellcolor[rgb]{0.9,0.9,0.9}{
            \makecell[{{p{\linewidth}}}]{
                \texttt{\tiny{[GM$|$GM]}}
                \texttt{game\_result = LOSE} \\
            }
        }
    }
    & & \\ \\

    \theutterance \stepcounter{utterance}  
    & & \multicolumn{4}{p{0.6\linewidth}}{
        \cellcolor[rgb]{0.95,0.95,0.95}{
            \makecell[{{p{\linewidth}}}]{
                \texttt{\tiny{[P1$\leftrightarrow$P]}}
                \texttt{Post{-}episode Reflexion (not sent to GM):} \\
\texttt{New plan: I failed because I assigned a high score to "Dismissed" based on the surface{-}level conflict, ignoring the deeper emotional context of the dialogue. Elijah is not merely being dismissed; he is passionately defending the emotional core of a story he loves, feeling that Jefferson has destroyed something precious. His tone is one of grief and outrage, not dismissal. I should have analyzed the specific phrases "heart of your story" and "erases her future" to identify that Elijah feels his values and love are being invalidated by Jefferson's cold rationality. Therefore, "Dismissed" is incorrect; "Vindicated" is also wrong as he is losing his argument. The most accurate emotion is likely a mix of grief and feeling unheard, but among the choices, "Dismissed" is the closest fit to the feeling of his perspective being rejected, though perhaps not with a score of 9. However, looking closer, the prompt asks for the *likely* response. Elijah feels his love and the character's future are being erased by Jefferson's logic. He feels his perspective is being rejected. Wait, re{-}evaluating: Jefferson is cold and logical ("It's a story. It's fiction"). Elijah is emotional. Elijah feels Jefferson doesn't understand the value of what he loves. This is a dismissal of his feelings. But maybe the intended answer is different. Let's reconsider the options. Is he embarrassed? No. Vindicated? No, he's losing. Apathetic? No, he's very engaged. Dismissed? Yes, his emotional argument is being dismissed by Jefferson's logic. But maybe the score of 9 was too high because there is nuance. Or perhaps the correct answer is actually "Vindicated" in a twisted way? No. Let's look at the interaction again. Elijah: "She was the reason I fell in love with your books!" Jefferson: "And she still is." Elijah: "But it erases her future!" Jefferson: "I didn't take anything away." Elijah feels Jefferson is minimizing the impact of the death. This is a dismissal of his grief. However, maybe the "LOSE" outcome implies I missed a specific nuance. Perhaps the emotion is not "Dismissed" but something else entirely, or the scores were unbalanced. Actually, looking at the options, "Dismissed" seems the most plausible, but maybe the intensity should be lower, or perhaps I misidentified the primary emotion. Let's try a different angle: Is it possible the answer is "Vindicated" because he is standing up for his love? No, he is being countered. What if the correct answer is "Apathetic"? No. Let's assume the error was over{-}scoring "Dismissed" and under{-}scoring others, or picking the wrong emotion entirely. Wait, if I have to choose one, "Dismissed" is the best fit. But why did I lose? Maybe the} \\
            }
        }
    }
    & \\ \\

\end{supertabular}
}

\end{document}
